% Compile with XeLaTeX (twice) or Tectonic. UTF-8 source.
% Standalone document: no external images or source files are required.
\documentclass[12pt,a4paper,oneside]{article}
\usepackage{fontspec}
\setmainfont{Times New Roman}
\setmonofont{Courier New}
\usepackage[english,russian]{babel}
\usepackage{geometry}
\geometry{left=2.54cm,right=2.54cm,top=2.54cm,bottom=2.54cm}
\usepackage{amsmath,array,tabularx,booktabs}
\usepackage{listings,xcolor,hyperref,enumitem,titlesec}
\titleformat*{\section}{\large\bfseries}
\titleformat*{\subsection}{\normalsize\bfseries}
\setcounter{secnumdepth}{0}
\hypersetup{pdftitle={Отчёт по ЛР № 5. ОС. Вариант 34},
  pdfauthor={А. Е. Пронина},hidelinks}
\setlength{\parindent}{0pt}
\setlength{\parskip}{4pt}
\setlist{nosep,leftmargin=1.5em}
\sloppy
\lstset{basicstyle=\fontsize{9}{10.5}\selectfont\ttfamily,frame=single,breaklines=true,
  columns=fullflexible,numbers=left,numberstyle=\tiny,
  showstringspaces=false,tabsize=4,keepspaces=true,
  aboveskip=7pt,belowskip=7pt}
\newcommand{\code}[1]{\texttt{\detokenize{#1}}}
\newcommand{\newsection}[1]{\clearpage\section{#1}}
\newcolumntype{Y}{>{\raggedright\arraybackslash}X}
\begin{document}
\begin{titlepage}
\begin{center}
\large МОСКОВСКИЙ АВИАЦИОННЫЙ ИНСТИТУТ\\
(НАЦИОНАЛЬНЫЙ ИССЛЕДОВАТЕЛЬСКИЙ УНИВЕРСИТЕТ)
\vspace{20pt}

Институт №8 «Компьютерные науки и прикладная математика»\\
Кафедра 806 «Вычислительная математика и программирование»
\end{center}
\vspace{65pt}
\begin{center}
\large\bfseries Отчет по лабораторной работе №5\\[7pt]
по курсу «Операционные системы»
\end{center}
\vspace{40pt}
\begin{center}
\large\bfseries Тема: «Динамические библиотеки»\\[7pt]
Вариант 34
\end{center}
\vfill
\begin{flushright}
\begin{tabular}{rl}
Выполнила: & А.\,Е. Пронина \\
Группа: & М8О-203БВ-25 \\
Преподаватель: & Е.\,С. Миронов \\
\end{tabular}
\end{flushright}
\vspace{40pt}
\begin{center}Москва, 2026\end{center}
\end{titlepage}
\setcounter{page}{2}

\section{Цель работы}
Приобрести практические навыки создания динамических библиотек и их
использования двумя способами: с заданием зависимости на этапе линковки
и с явной загрузкой во время выполнения программы. Изучить общий
интерфейс функций, получение их адресов и смену реализации без перезапуска.

\section{Постановка задачи}
Варианту 34 соответствуют функции №7 и №8 из задания.
Требуется создать две динамические библиотеки с одинаковыми контрактами
и разными реализациями.

\begin{center}
\small
\begin{tabularx}{\linewidth}{@{}p{0.30\linewidth}YY@{}}
\toprule
Контракт & Реализация 1 & Реализация 2\\
\midrule
\code{float Square(float A, float B)} &
Площадь прямоугольника: $S=AB$ &
Площадь прямоугольного треугольника: $S=AB/2$\\[5pt]
\code{char* translation(long x)} &
Двоичная запись числа & Троичная запись числа\\
\bottomrule
\end{tabularx}
\end{center}

Программа №1 должна использовать одну из библиотек, выбранную при
линковке. Программа №2 должна загружать библиотеки по относительным
путям, получать адреса функций и переключать реализации по команде
\code{0}. Обе программы предоставляют доступ к обеим функциям:
\begin{itemize}
\item \code{1 A B} --- вычислить площадь;
\item \code{2 x} --- перевести целое число;
\item \code{0} --- сменить библиотеку, только в программе №2;
\item конец ввода (EOF) --- завершить работу.
\end{itemize}

Проект реализован на C++11 с CMake для Linux и Windows. Используются
стандартные заголовки \code{iostream}, \code{cmath}, \code{climits}
и системный заголовок \code{dlfcn.h} либо \code{windows.h}.

\section{Основные теоретические сведения}
Динамическая библиотека хранит исполняемый код в отдельном файле:
\code{.so} в Linux или \code{.dll} в Windows. При задании зависимости
на этапе линковки файл библиотеки всё равно нужен во время запуска:
код библиотеки не включается в программу как при статической линковке.

При явной загрузке программа самостоятельно открывает библиотеку,
находит функции по именам и вызывает их через указатели. Общий
заголовок задаёт API. Для бинарной совместимости также должны совпадать
имена символов, типы, соглашение вызова и правила владения памятью.

\newsection{Архитектура проекта и общий интерфейс}
\begin{center}
\small
\begin{tabularx}{\linewidth}{@{}p{0.30\linewidth}Y@{}}
\toprule
Файл & Назначение\\
\midrule
\code{include/api.h} & Контракты, экспорт функций, типы указателей.\\
\code{src/impl1.cpp} & Прямоугольник и двоичная система.\\
\code{src/impl2.cpp} & Треугольник и троичная система.\\
\code{src/program1.cpp} & Обычные вызовы функций выбранной библиотеки.\\
\code{src/program2.cpp} & Загрузка, поиск функций, смена и освобождение библиотеки.\\
\code{src/console.h} & Обработка команд программы №2 через указатели.\\
\code{CMakeLists.txt} & Сборка двух библиотек и двух программ.\\
\bottomrule
\end{tabularx}
\end{center}

Обе программы могут использовать один файл \code{impl1}, но работают
в отдельных процессах. Команда \code{0} изменяет подключение только
программы №2. Она одновременно выбирает обе функции из новой библиотеки;
файлы библиотек на диске при этом не изменяются.

\subsection{Экспорт и соглашение вызова}
Ниже приведена платформенная часть \code{api.h}.
\begin{lstlisting}[language=C++]
#ifdef _WIN32
#define LAB5_CALL __cdecl
#if defined(LAB5_EXPORTS)
#define LAB5_API __declspec(dllexport)
#elif defined(LAB5_IMPORTS)
#define LAB5_API __declspec(dllimport)
#else
#define LAB5_API
#endif
#else
#define LAB5_CALL
#define LAB5_API
#endif
\end{lstlisting}
CMake определяет \code{LAB5_EXPORTS} при сборке DLL и
\code{LAB5_IMPORTS} для программы №1. Для Linux эти макросы
\code{LAB5_API} и \code{LAB5_CALL} заменяются пустым текстом.

\subsection{Контракты функций}
В C++ объявления помещены в блок \code{extern "C"}, чтобы экспортировать
имена \code{Square} и \code{translation} без C++ декорирования.
\begin{lstlisting}[language=C++]
extern "C" {
LAB5_API float LAB5_CALL Square(float A, float B);
LAB5_API char* LAB5_CALL translation(long x);
}
typedef float (LAB5_CALL *SquareFunction)(float, float);
typedef char* (LAB5_CALL *TranslationFunction)(long);
\end{lstlisting}
Типы указателей нужны программе №2.

\newsection{Реализации функций библиотек}
\subsection{Вычисление площади}
В первой библиотеке функция возвращает произведение сторон:
\begin{lstlisting}[language=C++]
float LAB5_CALL Square(float A, float B) {
    return A * B;
}
\end{lstlisting}
Во второй библиотеке стороны рассматриваются как катеты:
\begin{lstlisting}[language=C++]
float LAB5_CALL Square(float A, float B) {
    return A * B / 2.0f;
}
\end{lstlisting}
Проверка конечности и неотрицательности сторон выполняется в программах.
Нулевая сторона допустима. Точность результата ограничена типом
\code{float}; конечность входных данных не исключает переполнение произведения.

\subsection{Перевод в двоичную систему}
Основная часть \code{translation} из \code{impl1.cpp}:
\begin{lstlisting}[language=C++]
char* LAB5_CALL translation(long x) {
    static char buffer[sizeof(long) * CHAR_BIT + 2];
    int position = sizeof(buffer) - 1;
    buffer[position] = '\0';
    unsigned long number = static_cast<unsigned long>(x);
    if (x < 0) number = 0 - number;
    do {
        buffer[--position] = '0' + number % 2;
        number /= 2;
    } while (number != 0);
    if (x < 0) {
        buffer[--position] = '-';
    }
    return buffer + position;
}
\end{lstlisting}
Остатки от деления записываются справа налево. В \code{impl2.cpp}
вместо \verb|% 2| и \code{/= 2} используются \verb|% 3| и \code{/= 3}.
Цикл \code{do ... while} обеспечивает результат \code{"0"} для нуля.
Беззнаковая арифметика корректно обрабатывает модуль \code{LONG_MIN}.

Если \code{long} занимает $N$ бит, достаточно не более $N$ двоичных
цифр, одного минуса и завершающего нуля. Поэтому массив из $N+2$
символов не переполняется; для троичной системы нужно не больше места.
Возвращается адрес начала заполненной части массива.

Массив принадлежит библиотеке. Его нельзя освобождать и использовать
после выгрузки библиотеки; следующий вызов той же функции перезаписывает
результат. Решение рассчитано на однопоточную работу. Время перевода
пропорционально числу цифр, дополнительная память ограничена размером буфера.

\newsection{Два способа вызова функций}
\subsection{Программа №1: связь при линковке}
Зависимость задаётся в CMake через \code{target_link_libraries}.
В программе используются обычные вызовы, например:
\begin{lstlisting}[language=C++]
std::cout << Square(A, B) << '\n';
std::cout << translation(x) << '\n';
\end{lstlisting}
При запуске ОС загружает выбранную DLL/SO. По умолчанию используется
\code{impl1}; настройка \code{LINKED_IMPL=2} выбирает \code{impl2}
при новой сборке. Команда \code{0} в программе №1 не поддерживается.

\subsection{Программа №2: загрузка во время работы}
Для выбранной библиотеки хранятся дескриптор и адреса обеих функций:
\begin{lstlisting}[language=C++]
struct Library {
    LibraryHandle handle = nullptr;
    SquareFunction square = nullptr;
    TranslationFunction translate = nullptr;
};
\end{lstlisting}
В Windows тип \code{LibraryHandle} обозначает \code{HMODULE},
в Linux --- \code{void*}. Функция \code{open_library} открывает файл
и заполняет все три поля. Пути задаются относительно рабочего каталога:
\begin{lstlisting}[language=C++]
const char* paths[] = {"./" LAB5_LIBRARY1,
                       "./" LAB5_LIBRARY2};
\end{lstlisting}
Точные имена файлов передаются через CMake с помощью
\code{$<TARGET_FILE_NAME:impl1>} и аналогичного выражения для \code{impl2}.

\subsection{Загрузка в Linux}
Ключевые операции \code{open_library} приведены ниже; между ними в
исходнике проверяются результаты и освобождаются ресурсы при ошибке.
\begin{lstlisting}[language=C++]
library.handle = dlopen(path, RTLD_NOW | RTLD_LOCAL);
dlerror();
library.square = reinterpret_cast<SquareFunction>(
    dlsym(library.handle, "Square"));
const char* error = dlerror();
dlerror();
library.translate = reinterpret_cast<TranslationFunction>(
    dlsym(library.handle, "translation"));
error = dlerror();
\end{lstlisting}
\code{RTLD_NOW} требует разрешить необходимые символы до возврата из
\code{dlopen}; \code{RTLD_LOCAL} ограничивает видимость символов
для последующих загрузок. Перед \code{dlsym} сбрасывается прежняя ошибка,
после вызова проверяется \code{dlerror}. Тип указателя должен точно
соответствовать контракту функции.

\newsection{Загрузка в Windows и переключение реализации}
В Windows применяются другие функции ОС. Основные операции:
\begin{lstlisting}[language=C++]
library.handle = LoadLibraryA(path);
library.square = reinterpret_cast<SquareFunction>(
    GetProcAddress(library.handle, "Square"));
library.translate = reinterpret_cast<TranslationFunction>(
    GetProcAddress(library.handle, "translation"));
\end{lstlisting}
После каждого вызова проверяется результат. При ошибке сразу читается
\code{GetLastError}; если библиотека уже открыта, она освобождается.
В проекте используются относительные пути с ASCII-именами библиотек.

\subsection{Алгоритм команды 0}
При запуске \code{active=0}, то есть выбрана \code{impl1}. Другая
реализация определяется выражением \code{1 - active}. Основная часть
обработки команды из \code{program2.cpp}:
\begin{lstlisting}[language=C++]
int next = 1 - active;
Library replacement;
if (!open_library(paths[next], replacement)) {
    std::cerr
        << "Switch failed; current library remains active\n";
    continue;
}
Library previous = current;
current = replacement;
active = next;
if (!close_library(previous)) {
    result = 1;
    break;
}
std::cout << "Library " << active + 1 << '\n';
\end{lstlisting}
Сначала загружается новая библиотека и проверяются обе функции.
Только после успеха заменяются указатели и освобождается старый
дескриптор. Если новая библиотека или нужный символ не найдены,
текущая реализация остаётся доступной.

\subsection{Вызовы и освобождение ресурсов}
Функция \code{execute_command} получает \code{current.square} и
\code{current.translate}. Она читает аргументы и вызывает выбранную
функцию через указатель. Переключение затрагивает обе функции одновременно.

\code{close_library} вызывает \code{FreeLibrary} в Windows либо
\code{dlclose} в Linux, проверяет результат и обнуляет поля структуры.
После выгрузки прежние указатели не используются. При выходе
освобождается активная библиотека.

Ошибочные аргументы приводят к завершению с кодом 1. Неизвестная числовая
команда выводит сообщение и позволяет продолжить. EOF завершает работу
штатно; в Windows его можно передать сочетанием Ctrl+Z, затем Enter,
в Linux --- Ctrl+D на пустой строке.

\newsection{Организация сборки и запуск}
Основные цели CMake (часть настроек здесь опущена):
\begin{lstlisting}[language={}]
add_library(impl1 SHARED src/impl1.cpp)
add_library(impl2 SHARED src/impl2.cpp)
foreach(library impl1 impl2)
    target_include_directories(${library} PUBLIC include)
    target_compile_definitions(${library} PRIVATE LAB5_EXPORTS)
endforeach()

add_executable(program1 src/program1.cpp)
target_compile_definitions(program1 PRIVATE LAB5_IMPORTS)
target_link_libraries(program1 PRIVATE impl${LINKED_IMPL})

add_executable(program2 src/program2.cpp)
target_include_directories(program2 PRIVATE include)
add_dependencies(program2 impl1 impl2)
\end{lstlisting}
\code{SHARED} создаёт динамические библиотеки; в Unix CMake обеспечивает
позиционно независимый код. \code{add_dependencies} задаёт порядок сборки,
но не связывает программу №2 с реализациями. Для Linux/POSIX при необходимости
подключается библиотека загрузчика через \code{CMAKE_DL_LIBS}.

\subsection{Linux}
Команды выполняются из каталога \code{OS_lab5}:
\begin{lstlisting}[language=bash]
cmake -S . -B build -DCMAKE_BUILD_TYPE=Release
cmake --build build --parallel
cd build/bin
./program1
# Finish program1 before starting program2.
./program2
\end{lstlisting}
Для программы №1 задан путь поиска \code{$ORIGIN}, то есть каталог
исполняемого файла. Программа №2 использует относительные пути, поэтому
запускается из \code{build/bin}.


\newsection{Контрольные примеры и анализ способов подключения}
В проекте предусмотрены сборка и запуск на Linux, а также платформенная
ветка Windows. Запуск и тестирование на Windows не проводились.

\subsection{Вычисления и смена реализации}
В таблице приведены контрольные команды программы №2 и ожидаемые результаты.
Начальное меню опущено; сначала активна \code{impl1}.
\begin{center}
\begin{tabularx}{\linewidth}{@{}p{0.23\linewidth}p{0.25\linewidth}Y@{}}
\toprule
Ввод & Ожидаемый вывод & Пояснение\\
\midrule
\code{1 3 4} & \code{12} & Площадь прямоугольника.\\
\code{2 10} & \code{1010} & Двоичная запись.\\
\code{2 0} & \code{0} & Обработка нуля.\\
\code{2 -10} & \code{-1010} & Отрицательное число.\\
\code{0} & \code{Library 2} & Переход к \code{impl2}.\\
\code{1 3 4} & \code{6} & Площадь треугольника.\\
\code{2 10} & \code{101} & Троичная запись.\\
\code{2 -10} & \code{-101} & Отрицательное число.\\
\code{0} & \code{Library 1} & Возврат к \code{impl1}.\\
\code{1 3 4} & \code{12} & Снова прямоугольник.\\
\bottomrule
\end{tabularx}
\end{center}
Для варианта 34 результаты разных реализаций намеренно различаются:
меняются фигура и основание системы счисления.


\subsection{Сравнение способов подключения}
\begin{center}
\footnotesize
\begin{tabularx}{\linewidth}{@{}p{0.26\linewidth}YY@{}}
\toprule
Свойство & Программа №1 & Программа №2\\
\midrule
Выбор реализации & При линковке. & Во время выполнения.\\[2pt]
Вызов функции & По обычному объявлению из заголовка. & Через указатель,
полученный по имени символа.\\[2pt]
Загрузка файла & Автоматически загрузчиком ОС. & Явно через
\code{dlopen} или \code{LoadLibraryA}.\\[2pt]
Отсутствие библиотеки & Загрузчик не запускает программу при отсутствии
обязательной DLL/SO. & Обрабатывается программой; при неудачной смене
сохраняется старая реализация.\\[2pt]
Смена без перезапуска & Не предусмотрена. & Команда \code{0}.\\[2pt]
Сложность & Меньше служебного кода. & Нужны проверка символов,
хранение указателей и освобождение дескрипторов.\\
\bottomrule
\end{tabularx}
\end{center}

\section{Вывод}
Созданы две библиотеки с общим C ABI и две программы на C++.
Связь при линковке упрощает вызовы функций. Явная загрузка позволяет
сменить обе реализации без перезапуска, но требует проверки ошибок
и управления временем жизни библиотеки и указателей.

\end{document}
